\documentclass[../main.tex]{subfiles} % 使用 subfiles 文档类，指定主文档
\graphicspath{{\subfix{../image/}}} % 指定图片目录，后续可以直接使用图片文件名。

% 例如:
% \begin{figure}[H]
% \centering
% \includegraphics[width=0.3\textwidth]{图.png}
% \caption{}
% \label{figure:图}
% \end{figure}
% 注意:上述\label{}一定要放在\caption{}之后,否则引用图片序号会只会显示??.

\begin{document}
% SOURCE p=58; printed=43; object=ch02-title; region=top
\chapter{积分}\label{chapter:ch02-integration}

% SOURCE p=58; printed=43; object=ch02-intro; region=middle
在 $\mathbb{R}$ 上的经典积分理论中，$\int_a^b f(x)\,\mathrm{d}x$ 定义为 Riemann 和的极限；这些和是某些函数的积分，这些函数逼近 $f$，并且在 $[a,b]$ 的各个子区间上为常数。类似地，在任意测度空间上，对于在某种适当意义下局部常值的函数，都有显然的积分概念，并可将它扩展为更一般函数的积分。本章建立抽象测度空间上的积分理论，特别关注 $\mathbb{R}$ 上的 Lebesgue 积分及其向 $\mathbb{R}^n$ 的推广。
\subfile{Chapter-02/section-01.tex}
\subfile{Chapter-02/section-02.tex}
\subfile{Chapter-02/section-03.tex}
\subfile{Chapter-02/section-04.tex}
\subfile{Chapter-02/section-05.tex}
\subfile{Chapter-02/section-06.tex}
\subfile{Chapter-02/section-07.tex}
\subfile{Chapter-02/section-08.tex}
\end{document}
